%============================================================= §6 DEFENSE
\section{Countermeasure: Multi-Vector Active Defense Engine}\label{sec:defense}

Requirement~(d) asks whether a defensive countermeasure was developed to mitigate the attack. In
our project, this countermeasure constitutes our qualifying bonus deliverable: an autonomous,
host-based active defense monitor running on the victim node that continuously analyzes raw Layer-2
traffic, detects scanning patterns in real time, and dynamically severs the attacker's connection
via kernel firewall rule injection.

\subsection{Detection Methodology \& The Four Inspection Vectors}
The week-11 Design Report proposed a single sliding-window detector. The delivered engine expands
this foundation into a four-vector inspection engine designed to detect both aggressive sweeps and
stealthy reconnaissance:

\begin{table}[htbp]
\centering
\small
\renewcommand{\arraystretch}{1.25}
\caption{Delivered Active Defense Detection Vectors and Trigger Mechanisms}
\label{tab:defense_vectors}
\begin{tabularx}{\textwidth}{@{}L{3.6cm} L{4.0cm} X@{}}
\toprule
\hdrrow \thh{Detection Vector} & \thh{Operational Trigger Threshold} & \thh{Security Threat Addressed}\\
\midrule
Fast SYN Sliding Window & $C \ge 8$ SYNs across $U \ge 5$ distinct ports within $W_1 = 10.0$\,s. & Detects rapid, multi-port automated scans while suppressing false positives from busy legitimate services.\\
\addlinespace
Slow / Stealth Window & $C \ge 20$ SYNs within $W_2 = 60.0$\,s. & Catches ``low-and-slow'' scans that inject inter-packet delays ($\ge 2.0$\,s) to evade the 10-second window.\\
\addlinespace
Canary Honeyport Trap & Touching unassigned decoy port (Port 81). & Zero-delay instant quarantine; legitimate clients have no reason to access unadvertised honeyports.\\
\addlinespace
Invalid Flag Filter & SYN+FIN (\texttt{0x03}), NULL (\texttt{0x00}), or XMAS. & Catches Phase~3 OS-fingerprinting rule-breaker probes on the first observation; zero legitimate usage.\\
\bottomrule
\end{tabularx}
\end{table}

\subsection{Cross-Platform Dynamic Firewall Mitigation Architecture}
When an alert is tripped with the \texttt{--block} flag active, the engine dynamically injects an
inbound block rule into the host operating system's kernel packet filter:
\begin{itemize}
    \item \textbf{Windows Defender Firewall:} Dynamically injects an inbound quarantine rule:
    \begin{center}
    \small\texttt{netsh advfirewall firewall add rule name=CustomBlock\_IP dir=in action=block remoteip=IP}
    \end{center}
    \item \textbf{Linux Netfilter:} Dynamically injects a kernel packet drop rule:
    \begin{center}
    \small\texttt{iptables -I INPUT -s IP -p tcp -j DROP}
    \end{center}
\end{itemize}
To ensure robust, production-grade operation, three safeguards are enforced:
\begin{enumerate}
    \item \textbf{Idempotent Rule Injection:} Tracks quarantined IP addresses in memory; burst alerts
    during an ongoing scan never create duplicate firewall rules.
    \item \textbf{Self-Healing Auto-Release Daemon:} A dedicated background thread (\texttt{threading.Thread})
    monitors quarantine TTL expirations independently of incoming traffic. Rules are deleted cleanly
    after a configurable timeout (default 60\,s or 300\,s) or upon SIGINT termination.
    \item \textbf{Structured JSON SIEM Audit Trail:} Every detection, rule injection, and release event
    is logged to \texttt{defense\_incident\_log.json} with ISO timestamps, attacker IPs, trigger
    reasons, and rule identifiers.
\end{enumerate}

\subsection{Sliding-Window Detection and Mitigation Timeline}
Figure~\ref{fig:timing_defense} illustrates the timeline of a fast port scan being detected and
quarantined by the victim engine:

\begin{figure}[htbp]
\centering
\begin{tikzpicture}[font=\sffamily\small, scale=0.88, every node/.style={transform shape}]
    % Time Axis
    \draw[thick, -{Stealth[length=6pt]}, color=ink] (0, 0) -- (14.5, 0) node[right, font=\scriptsize\bfseries] {Time (s)};
    
    % Sliding Observation Window Indicator Bar
    \draw[accent, thick, |-|] (0.5, 2.45) -- (8.8, 2.45) 
        node[midway, above=1pt, font=\scriptsize\bfseries, color=accent] {Sliding Observation Window ($W_1 = 10.0\,\text{s}$)};

    % Staggered SYN Probes
    \draw[thick, color=accent] (0.5, 0) -- (0.5, 0.75); \fill[accent] (0.5, 0) circle (2pt);
    \node[draw=rule, fill=codebg, rounded corners=2pt, font=\tiny\bfseries] at (0.5, 0.95) {SYN :80};

    \draw[thick, color=accent] (1.8, 0) -- (1.8, 1.45); \fill[accent] (1.8, 0) circle (2pt);
    \node[draw=rule, fill=codebg, rounded corners=2pt, font=\tiny\bfseries] at (1.8, 1.65) {SYN :22};

    \draw[thick, color=accent] (3.1, 0) -- (3.1, 0.75); \fill[accent] (3.1, 0) circle (2pt);
    \node[draw=rule, fill=codebg, rounded corners=2pt, font=\tiny\bfseries] at (3.1, 0.95) {SYN :21};

    \draw[thick, color=accent] (4.4, 0) -- (4.4, 1.45); \fill[accent] (4.4, 0) circle (2pt);
    \node[draw=rule, fill=codebg, rounded corners=2pt, font=\tiny\bfseries] at (4.4, 1.65) {SYN :443};

    \draw[thick, color=accent] (5.7, 0) -- (5.7, 0.75); \fill[accent] (5.7, 0) circle (2pt);
    \node[draw=rule, fill=codebg, rounded corners=2pt, font=\tiny\bfseries] at (5.7, 0.95) {SYN :445};

    \draw[thick, color=accent] (7.0, 0) -- (7.0, 1.45); \fill[accent] (7.0, 0) circle (2pt);
    \node[draw=rule, fill=codebg, rounded corners=2pt, font=\tiny\bfseries] at (7.0, 1.65) {SYN :135};

    \draw[thick, color=accent] (8.3, 0) -- (8.3, 0.75); \fill[accent] (8.3, 0) circle (2pt);
    \node[draw=rule, fill=codebg, rounded corners=2pt, font=\tiny\bfseries] at (8.3, 0.95) {SYN :139};

    % 8th Probe - THRESHOLD BREACH
    \draw[thick, color=accent] (9.6, 0) -- (9.6, 1.45); \fill[accent] (9.6, 0) circle (2pt);
    \node[draw=rule, fill=codebg, rounded corners=2pt, font=\tiny\bfseries] at (9.6, 1.65) {SYN :3389};

    % Time Labels
    \foreach \x/\timeval in {0.5/0.0, 1.8/1.0, 3.1/2.0, 4.4/3.0, 5.7/4.0, 7.0/5.0, 8.3/6.0, 9.6/7.0, 12.5/8.5} {
        \draw[thick] (\x, 0.08) -- (\x, -0.08) node[below=2pt, font=\tiny, fill=white, inner sep=1pt] {$t=\timeval$};
    }

    % Threshold Line
    \draw[thick, dashed, color=accentstrong] (9.6, 1.95) -- (9.6, 2.65);
    \node[draw=accentstrong, fill=accentsoft, rounded corners=2pt, font=\tiny\bfseries, align=center] at (9.6, 3.05) {
        \color{accentstrong}THRESHOLD BREACHED\\
        \tiny $C = 8 \ge \theta_{\text{count}}, \ U = 8 \ge \theta_{\text{ports}}$
    };

    % Firewall Injection Callout
    \draw[-{Stealth[length=4pt]}, thick, color=accentstrong] (9.6, -0.45) -- (9.6, -0.85);
    \node[draw=rule, fill=codebg, rounded corners=2pt, font=\tiny, align=center] at (9.6, -1.25) {
        \textbf{\color{accent}Dynamic Firewall Rule Injected via \texttt{netsh}}\\
        \texttt{\scriptsize netsh advfirewall add rule dir=in action=block remoteip=192.168.0.105}
    };

    % 9th Probe Blocked
    \draw[thick, color=accentstrong] (12.5, 0) -- (12.5, 1.3);
    \fill[accentstrong] (12.5, 0) circle (2pt);
    \node[draw=accentstrong, fill=accentsoft, rounded corners=2pt, font=\tiny\bfseries, align=center] at (12.5, 1.6) {
        SYN $\to$ :8080\\
        \color{accentstrong}[BLOCKED / DROPPED]
    };
    \draw[line width=1.1pt, color=accentstrong] (12.3, 0.4) -- (12.7, 0.8);
    \draw[line width=1.1pt, color=accentstrong] (12.3, 0.8) -- (12.7, 0.4);
\end{tikzpicture}
\caption{Timeline of Sliding-Window Scan Detection and Dynamic Firewall Rule Injection.}
\label{fig:timing_defense}
\end{figure}

\FloatBarrier
\subsection{Empirical Diagnostic Logs: Active Defense in Action}

\subsubsection{Evaluation 1: Windows Defender Firewall Dynamic Quarantine \& Decoy Honeyport}
The log below captures the Windows defense agent detecting a fast scan, tripping on decoy port~81,
trapping an invalid SYN+FIN probe, and releasing quarantine via the background TTL daemon:

\begin{logbox}
=============================================================================
DEFENSIVE MONITOR & INCIDENT MITIGATION TELEMETRY
Platform: Windows 11 Pro Desktop (192.168.0.116) | Adapter: Npcap Loop
Mitigation: Active Dynamic Windows Defender Firewall Rule Injection (--block)
=============================================================================

[SCENARIO A: Fast SYN Scan Rate Detection]
  SYN 192.168.0.105 -> 192.168.0.116:8080  window_count=8  distinct_ports=8
  [ALERT] 192.168.0.105 | fast port scan (SYN=8, distinct ports=8 within 10.0s)
  [MITIGATION] Blocked TCP traffic from 192.168.0.105
               Firewall rule: CustomPortScanBlock_192_168_0_105 (Windows Defender Firewall)
               Status: Dynamic inbound blocking rule active across all profiles.

[SCENARIO B: Canary Honeyport Decoy Port 81 Trapping]
  SYN 192.168.0.105 -> 192.168.0.116:81
  [ALERT] 192.168.0.105 | hit canary honeyport (81)  [Zero-Delay Instant Quarantine]
  [MITIGATION] Injected quarantine rule: CustomPortScanBlock_192_168_0_105 (TTL: 60s)

[SCENARIO D: RFC 793 Rule-Breaker Flag Anomaly Trapping]
  SYN+FIN 192.168.0.105 -> 192.168.0.116:21
  [CRITICAL ALERT] invalid TCP flag combination (SYN+FIN, 0x03) on port 21
  [MITIGATION] Severity Critical: Injected quarantine rule (Quarantine TTL: 300s)

[SCENARIO F: Self-Healing Auto-Release Background Daemon]
  [MITIGATION] Released quarantine block on 192.168.0.105 (TTL 60s elapsed)
  Status: Invoked `netsh advfirewall firewall delete rule name=CustomPortScanBlock_192_168_0_105`
          Rule cleanly deleted from kernel table; client reachability restored.
\end{logbox}

\subsubsection{Evaluation 2: Linux Netfilter Dynamic Rule Injection \& Drop Counter Progression}
On Linux, the defense agent injects native Netfilter rules. The log verifies the kernel DROP counter
climbing dynamically as probe packets are dropped before reaching application sockets:

\begin{logbox}
=============================================================================
LINUX NETFILTER ACTIVE MITIGATION TELEMETRY
Platform: Kali Linux Kernel 6.x VM (192.168.56.101) | Interface: eth0
=============================================================================
[ALERT] 192.168.56.1 | fast port scan (SYN=8, distinct ports=8 within 10.0s)
[MITIGATION] Injected Linux Netfilter rule: CustomPortScanBlock_192_168_56_1
Command: `sudo iptables -I INPUT -s 192.168.56.1 -p tcp -j DROP`

Kernel Packet Drop Counter Audit (`sudo iptables -L INPUT -n -v`):
  Chain INPUT (policy ACCEPT 418 packets)
  pkts bytes target     prot opt in     out     source          destination
    12   720 DROP       tcp  --  *      *       192.168.56.1    0.0.0.0/0
    13   780 DROP       tcp  --  *      *       192.168.56.1    0.0.0.0/0  (Subsequent probe dropped)
    14   840 DROP       tcp  --  *      *       192.168.56.1    0.0.0.0/0  (Subsequent probe dropped)
    16   960 DROP       tcp  --  *      *       192.168.56.1    0.0.0.0/0  (Attacker severed completely)
\end{logbox}

\subsubsection{Evaluation 3: Live Mid-Attack Chain Severance}
The definitive operational proof of the active defense bonus is its ability to sever a live attack
chain mid-execution. During our live end-to-end verification run, the scanner completed its first
exploit (FTP Anonymous login), but while attempting subsequent exploitation modules, the defense
engine's firewall rule took effect. All subsequent exploits timed out with \texttt{WinError 10060}:

\begin{logbox}
[VICTIM DEFENSE AGENT LOG]
  SYN 192.168.0.105 -> 192.168.0.116:55455  window_count=10  distinct_ports=5
  [ALERT] Possible TCP port scan from 192.168.0.105
  [MITIGATION] Blocked TCP traffic from 192.168.0.105
               Firewall rule: CustomPortScanBlock_192_168_0_105

[ATTACKER WORKSTATION CONSOLE - ATTACK CHAIN SEVERED MID-EXECUTION]
  [EXPLOIT 1] FTP Anonymous Login (port 21) -> SUCCESS [OK] (Executed before threshold breach)
  [EXPLOIT 2] FTP Write -> Web Shell -> RCE
    Status : FAILED [X]
    Detail : Connection failed: timed out ([WinError 10060] A connection attempt failed because
             the connected party did not properly respond after a period of time)
  [EXPLOIT 3] FTP Bounce Attack Pivot -> FAILED [X] (Connection timed out)
  [EXPLOIT 4] HTTP Directory Traversal -> FAILED [X] (All 4 traversal payloads timed out)
\end{logbox}
